
\chapter{三字經}

\begin{verse}
    \vinsm{一} 人之初性本善性相近習相遠
    \vinsm{二} 苟不教性乃遷教之道貴以專
    \vinsm{三} 昔孟母擇鄰處子不學斷機杼
    \vinsm{四} 竇燕山有義方教五子名俱揚
    \vinsm{五} 養不教父之過教不嚴師之惰
    \vinsm{六} 子不學非所宜幼不學老何為
    \vinsm{七} 玉不琢不成器人不學不知義
    \vinsm{八} 為人子方少時親師友習禮儀
    \vinsm{九} 香九齢能溫席孝於親所當執
    \vinsm{十} 融四歳能讓梨弟於長宜先知
    \vinsm{十一} 首孝弟次見聞知某數識某文
    \vinsm{十二} 一而十十而百百而千千而萬
    \vinsm{十三} 三才者天地人三光者日月星
    \vinsm{十四} 三綱者君臣義父子親夫婦順
    \vinsm{十五} 曰春夏曰秋冬此四時運不窮
    \vinsm{十六} 曰南北曰西東此四方應乎中
    \vinsm{十七} 曰水火木金土此五行本乎數
    \vinsm{十八} 曰仁義禮知信此五常不容紊
    \vinsm{十九} 稻粱菽麥黍稷此六穀人所食
    \vinsm{二十} 馬牛羊雞犬豕此六畜人所飼
    \vinsm{二十一} 曰喜怒曰哀懼愛惡欲七情具
    \vinsm{二十二} 匏土革木石金與絲竹乃八音
    \vinsm{二十三} 高曾祖父而身身而子子而孫自子孫至\footnote{原作曰至元曾避諱}玄曾乃九族人之倫
    \vinsm{二十四} 父子恩夫婦從兄則友弟則恭長則敬臣則忠此十義人所同
    \vinsm{二十五} 凡訓蒙須講究詳訓詁名句讀
    \vinsm{二十六} 為學者必有初小學終至四書
    \vinsm{二十七} 論語者二十篇群弟子記善言
    \vinsm{二十八} 孟子者七篇止講道德說仁義
    \vinsm{二十九} 作中庸子思筆中不偏庸不易
    \vinsm{三十} 作大學乃曾子自修齊至平治
    \vinsm{三十一} 孝經通四書熟如六經始可讀
    \vinsm{三十二} 詩書易禮春秋號六經當講求
    \vinsm{三十三} 有連山有歸藏有周易三易詳
    \vinsm{三十四} 有典謨有訓誥有誓命書之奧
    \vinsm{三十五} 我周公作周禮著六官存治體
    \vinsm{三十六} 大小戴註禮記述聖言禮樂備
    \vinsm{三十七} 曰國風
    \vinsm{三十八}
    \vinsm{三十九}
    \vinsm{四十}
    \vinsm{四十一}
    \vinsm{四十二}
    \vinsm{四十三}
    \vinsm{四十四}
    \vinsm{四十五}
    \vinsm{四十六}
    \vinsm{四十七}
    \vinsm{四十八}
    \vinsm{四十九}
    \vinsm{五十}
    \vinsm{五十一}
    \vinsm{五十二}
    \vinsm{五十三}
    \vinsm{五十四}
    \vinsm{五十五}
    \vinsm{五十六}
    \vinsm{五十七}
    \vinsm{五十八}
    \vinsm{五十九}
    \vinsm{六十}
    \vinsm{六十一}
    \vinsm{六十二}
    \vinsm{六十三}
    \vinsm{六十四}
    \vinsm{六十五}
    \vinsm{六十六}
    \vinsm{六十七}
    \vinsm{六十八}
    \vinsm{六十九}
    \vinsm{七十}
    \vinsm{七十一}
    \vinsm{七十二}
    \vinsm{七十三}
    \vinsm{七十四}
    \vinsm{七十五}
    \vinsm{七十六}
    \vinsm{七十七}
    \vinsm{七十八}
    \vinsm{七十九}
    \vinsm{八十}
    \vinsm{八十一}
    \vinsm{八十二}
    \vinsm{八十三}
    \vinsm{八十四}
    \vinsm{八十五}
    \vinsm{八十六}
\end{verse}
